\documentclass[11pt]{article}
\usepackage{acl}
\usepackage{amsmath,amssymb}
\usepackage{graphicx}
\usepackage{booktabs}
\usepackage[colorlinks=false,pdfborder={0 0 0}]{hyperref}

% ACL-style version: adapted framing for ML/NLP audience
% Emphasizes methodology and decision-theoretic aspects

\title{From Prediction to Decision: Uncertainty-Aware Deferral for Reliable Medical Image Segmentation}

\author{Author Names Withheld}

\begin{document}

\maketitle

\begin{abstract}
Dense prediction models (segmentation, tagging, parsing) produce per-element
outputs, but deployment requires knowing which outputs to trust. We study this
as selective prediction for dense labeling, using medical image segmentation as
a testbed. We compare two uncertainty methods (Monte Carlo Dropout, Test-Time
Augmentation) under three deferral policies (global threshold, image-adaptive,
confidence-aware) and evaluate temperature scaling as a calibration step. On
retinal vessel segmentation across three datasets, we find: (1) TTA uncertainty
achieves 0.881 AUROC for error prediction versus 0.768 for MC Dropout at
6.3$\times$ lower latency; (2) a confidence-aware deferral score $s = u \cdot
(1 - c)$, where $u$ is uncertainty and $c$ is prediction confidence, reduces
error by 55\% at 12\% deferral rate (4.6$\times$ more efficient than global
thresholding); (3) temperature scaling does not improve deferral quality; (4)
uncertainty AUROC transfers robustly under domain shift ($>$0.83 on external
datasets despite 5\% accuracy degradation). These findings generalize to any
dense prediction task where selective abstention is practical.
\end{abstract}

\section{Introduction}

Selective prediction, the idea that a model should abstain when uncertain, is
well studied for classification \citep{geifman2017selective,
el2010foundations}. Dense prediction tasks (segmentation, sequence labeling,
structured prediction) introduce additional complications: the unit of
abstention is unclear (pixel, region, sequence, image), uncertainty maps are
high-dimensional, and the interaction between uncertainty estimation and
deferral policy is poorly understood.

We use medical image segmentation as a testbed because the stakes are high
(clinical decisions depend on outputs), the evaluation is well defined (pixel
ground truth is available), and uncertainty methods are established (MC Dropout,
TTA, ensembles). The findings, however, apply more broadly.

Our question is not ``how do we estimate better uncertainty?'' but
``given an uncertainty estimate, how should we use it?'' This shifts
evaluation from calibration metrics (ECE) to decision metrics (error reduction
at a given deferral budget).

\section{Methods}

\paragraph{Uncertainty estimation.}
MC Dropout \citep{gal2016dropout}: $T = 30$ passes, mutual information as
uncertainty. TTA: $K = 6$ geometric transforms, variance as uncertainty.

\paragraph{Deferral policies.}
\textbf{Global:} defer if $u_{ij} > \tau$. \textbf{Adaptive:} per-instance
percentile threshold. \textbf{Confidence-aware:} score $s_{ij} = u_{ij} \cdot
(1 - c_{ij})$ with $c_{ij} = 2|\hat{p}_{ij} - 0.5|$; defer pixels that are
both uncertain and near the decision boundary.

\paragraph{Calibration.}
Temperature scaling: $\hat{p}^{\text{cal}} = \sigma(\text{logit}(\hat{p})/T)$.

\paragraph{Evaluation.}
Error reduction ratio (ERR), deferral rate, AUROC of uncertainty as an error
classifier (Unc-AUROC), area under the coverage curve (AUCC), and
deferral efficiency (ERR / deferral rate).

\section{Experimental setup}

U-Net (ResNet-34 encoder) for retinal vessel segmentation. Training: DRIVE
(20 images, patches, 80 epochs). Evaluation: DRIVE (20 test), STARE (20),
CHASE\_DB1 (28). Cross-dataset evaluation is zero-shot.

\section{Results}

\paragraph{Uncertainty quality.}
TTA Unc-AUROC: 0.881. MC Dropout: 0.768 ($p < 0.001$). Segmentation quality
comparable (Dice 0.768 vs.\ 0.764, $p = 0.142$). TTA runs 6.3$\times$
faster.

\paragraph{Deferral efficiency.}

\begin{table}[t]
\centering
\small
\begin{tabular}{llccc}
\toprule
Unc. & Deferral & ERR & Def.\% & Eff. \\
\midrule
MC & Global & 26.4 & 33.4 & 0.79 \\
MC & Conf-aware & 48.9 & 10.9 & 4.49 \\
TTA & Global & 42.2 & 20.0 & 2.11 \\
TTA & Adaptive & 79.5 & 25.0 & 3.18 \\
TTA & Conf-aware & 55.2 & 12.1 & 4.56 \\
\bottomrule
\end{tabular}
\caption{Deferral results. Eff.\ = ERR / Def.\%.}
\end{table}

At a fixed 10\% deferral budget, TTA + confidence-aware achieves 43.2\% ERR.
MC Dropout + global: 14.7\%. The 3$\times$ gap is the main practical finding.

\paragraph{Calibration.}
Temperature scaling: ECE improvement $<$ 0.001, Unc-AUROC degradation 0.019
(MC Dropout). Calibration and deferral optimize different objectives.

\paragraph{Domain shift.}
STARE: Dice $-$5\%, Unc-AUROC $+$10\%. CHASE\_DB1: similar pattern. Uncertainty
quality is more portable than prediction quality.

\section{Analysis}

TTA probes input-level prediction fragility; MC Dropout probes weight-space
variance. For dense prediction, input-level fragility localizes more precisely
to error-prone structures.

Confidence-aware deferral recognizes that uncertainty $\neq$ error. The score
$s = u(1-c)$ filters pixels that are uncertain but far from the decision
boundary, which are likely correct despite their uncertainty.

Temperature scaling is a monotonic transform: it changes probability magnitudes
but not rankings. Deferral policies based on ranking are invariant to it.
The only effect is through the confidence score $c$, which is second-order.

\section{Discussion}

The core finding generalizes beyond segmentation. Any dense prediction task
with a meaningful unit of abstention (pixels, tokens, time steps) can benefit
from the same framework: estimate uncertainty, weight by prediction confidence,
defer the highest-risk elements. The confidence-aware score $s = u(1-c)$ is
applicable whenever predictions have a natural decision boundary.

The calibration finding also generalizes. ECE measures a property of the
probability distribution, not of the decision policy. Good calibration
neither implies nor is implied by good selective prediction. Practitioners
should evaluate uncertainty methods by the metrics that match their use case.

\section{Conclusion}

Uncertainty methods should be judged by what they enable, not what they
measure. TTA with confidence-aware deferral provides the best error reduction
per deferred element across our experiments. Calibration is orthogonal to
deferral quality.

\bibliography{../references}
\bibliographystyle{acl_natbib}

\end{document}
